\section{Internalizing the Harness into a Small Student}
\label{sec:distill}

Two things could carry the harness's gains: the grounded scaffolding, or behavior absorbed into the weights. This section asks how much is the latter---whether the loop's interaction quality can be \emph{internalized} so a small model emits high-quality artifacts in fewer passes. It is a complement to the thesis, not a substitute: the cheaper internalized proposer is still meant to run \emph{inside} the same grounded harness, so internalization and grounding compose. We distill judge-filtered teacher trajectories from the harness into an $8$B Qwen3-VL student via supervised fine-tuning.

\paragraph{Result.} The student recovers about half the teacher's single-sample capability on an out-of-distribution suite, and a second sample closes much of the remaining gap (\Cref{fig:distill})---at roughly $10\times$ lower deployment cost than running the harness over a frontier model. On a harder held-out suite of the teacher's own pre-curated failures it passes a substantial fraction (\Cref{tab:distill-headline}), and run back inside the harness it recovers still more.

\begin{figure}[t]
\centering
\begin{tikzpicture}
\begin{axis}[paperbar,width=0.6\linewidth,height=4.6cm,clip=false,
  ybar,bar width=26pt,ymin=0,ymax=1.2,
  ylabel={capability vs.\ teacher ($\times$)},
  symbolic x coords={mean@1,pass@2},xtick=data,enlarge x limits=0.6,
  x tick label style={font=\footnotesize,color=cMut},
  ymajorgrids,ytick={0,0.25,0.5,0.75,1.0},yticklabel={\pgfmathprintnumber{\tick}$\times$}]
  % teacher reference band spanning the full plot width
  \fill[black!5] ({rel axis cs:0,0}|-{axis cs:mean@1,1.0}) rectangle ({rel axis cs:1,0}|-{axis cs:mean@1,1.2});
  \draw[cMut,densely dashed,line width=0.7pt] ({rel axis cs:0,0}|-{axis cs:mean@1,1.0}) -- ({rel axis cs:1,0}|-{axis cs:mean@1,1.0});
  \node[font=\footnotesize,cMut,anchor=south east] at ({rel axis cs:0.99,0}|-{axis cs:mean@1,1.01}) {teacher ($1\times$)};
  \addplot[draw=none,fill=cStu!65] coordinates {(mean@1,0.51) (pass@2,0.70)};
  \node[font=\footnotesize\bfseries,cInk,anchor=south] at (axis cs:mean@1,0.52) {0.51$\times$};
  \node[font=\footnotesize\bfseries,cInk,anchor=south] at (axis cs:pass@2,0.71) {0.70$\times$};
  \node[font=\scriptsize,cGnd!80!black,anchor=center]
    at ($(axis cs:mean@1,0.90)!0.5!(axis cs:pass@2,0.90)$)
    {sampling adds $+0.19\times$};
\end{axis}
\end{tikzpicture}
\caption{\textbf{An $8$B student internalizes much of the teacher's interaction quality.} Distilled student as a fraction of the frontier teacher on the $44$-task out-of-distribution suite (\Cref{tab:distill-ood}): one sample recovers $0.51\times$ and two samples $0.70\times$ of teacher capability, at ${\sim}10\times$ lower deployment cost. On the $18$-task hard held-out suite the same student passes $44\%$ at pass@$1$ and $56\%$ at pass@$3$ (\Cref{tab:distill-headline}).}
\label{fig:distill}
\end{figure}

\paragraph{Variance is a budget (a cautionary finding).} Reinforcement fine-tuning (RFT) on top of SFT improves every per-turn consistency metric we measure \emph{and} regresses pass@$k$ (\Cref{fig:variance}). The two students start nearly tied on a single sample; as samples are added, the SFT student keeps converting draws into new solved tasks while the RFT student's curve goes flat. The reading is that for sampling-based deployment, output-distribution variance \emph{is} the inference-scaling lever, and any post-training that reduces it spends from the same budget best-of-$N$ draws on. Variance is not noise to minimize; it is the resource sampling converts into quality. The operational recipe is therefore: SFT to internalize interaction quality, keep the proposer's sampling temperature, and run it inside the grounded harness.

\begin{figure}[t]
\centering
\begin{tikzpicture}
\begin{axis}[width=0.56\linewidth,height=4.6cm,
  xlabel={samples $k$},ylabel={held-out judge-keep (\%)},
  xmin=0.8,xmax=3.6,ymin=20,ymax=62,clip=false,
  xtick={1,2,3},xticklabels={mean@1,pass@2,pass@3},
  x tick label style={font=\footnotesize,color=cMut},
  ytick={25,35,45,55},ymajorgrids,
  axis x line*=bottom,axis y line*=left]
  \addplot[cStu,line width=1.4pt,mark=*,mark size=1.9,
    mark options={fill=cStu,draw=white,line width=0.5pt}]
    coordinates {(1,31)(2,50)(3,56)};
  \addplot[cUng,line width=1.1pt,mark=square*,mark size=1.7,
    mark options={fill=cUng,draw=white,line width=0.5pt}]
    coordinates {(1,28)(2,39)(3,39)};
  \node[font=\scriptsize\bfseries,cStu,anchor=west,inner sep=3pt] at (axis cs:3,56) {SFT student: 56};
  \node[font=\scriptsize,cUng!85!black,anchor=west,inner sep=3pt] at (axis cs:3,39) {+\,RFT: 39};
\end{axis}
\end{tikzpicture}
\caption{\textbf{RFT polishes consistency but spends the variance that sampling needs.} Judge-keep vs.\ number of samples on the $18$-task hard held-out suite for the chosen SFT student and the same student after RFT. RFT improves every per-turn consistency metric (\Cref{tab:distill-rft}) yet regresses pass@$3$ by $17$\,pp: the near-tie at one sample and the flat curve past it show the lost quality is exactly the distributional spread best-of-$N$ converts into solved tasks.}
\label{fig:variance}
\end{figure}
